\documentclass[11pt,a4paper]{article}

\usepackage[T1]{fontenc}
\usepackage[utf8]{inputenc}
\usepackage{lmodern}
\usepackage[a4paper,margin=26mm]{geometry}
\usepackage{amsmath,amssymb,amsthm,mathtools}
\usepackage{booktabs,array}
\usepackage{microtype}
\usepackage{xcolor}
\usepackage{enumitem}
\usepackage{url}
\usepackage[colorlinks=true,linkcolor=blue!55!black,citecolor=blue!55!black,
            urlcolor=blue!55!black]{hyperref}

\newtheorem{theorem}{Theorem}
\newtheorem{lemma}[theorem]{Lemma}
\newtheorem{proposition}[theorem]{Proposition}
\newtheorem{remark}[theorem]{Remark}
\newcommand{\Z}{\mathbb{Z}}
\newcommand{\R}{\mathcal{R}}
\newcommand{\E}{\mathbb{E}}
\newcommand{\Prb}{\mathop{\rm Pr}}
\newcommand{\Comp}{\mathsf{Comp}}
\newcommand{\Sign}{\mathsf{Sign}}
\newcommand{\Verify}{\mathsf{Verify}}
\newcommand{\KeyGen}{\mathsf{KeyGen}}
\newcommand{\pk}{\mathsf{pk}}
\newcommand{\sk}{\mathsf{sk}}
\newcommand{\norm}[1]{\left\lVert #1\right\rVert}

\title{\bfseries Branch-Dependent Compression Leakage in DARTS-128:\\
Ternary-Secret Recovery and a Fresh-Message Forgery\\
from 17 Million Signatures}
\author{Martin Feussner\\
\small Selmer Center, University of Bergen\\
\small \texttt{martin.feussner@uib.no}}
\date{26 September 2026}

\begin{document}
\maketitle

\begin{center}
\fbox{\parbox{0.92\linewidth}{\small\textbf{Provenance and review status.}
The attack was independently found by an AI system, OpenAI Codex using
Daybreak Blue at maximum reasoning effort, without a human first proposing
this specific attack.  The derivation, implementation, experiments, and
initial manuscript were also produced by the AI.  At the date above, no human
cryptanalyst has independently verified the argument, experiments, or novelty
assessment.  This is a preliminary disclosure for human review and
reproduction.}}
\end{center}
\vspace{0.7em}

\begin{abstract}
DARTS is a module-lattice signature submitted to the Next-generation
Commercial Cryptographic Algorithms Program.  Its signer forms
$z=y+(-1)^bcs$ with an initially balanced bimodal bit $b$.  Before the main
sphere rejection, Algorithm~7 accepts the commitment only if
$\Comp(w)=\Comp(w+(1-b)c\alpha)$.  This is a real filter for $b=0$ and the
identity for $b=1$.  A gate-only independence heuristic predicts a
branch-zero fraction of $0.44133$ among gate-passing attempts for DARTS-128.
The fully accepted fraction was $0.44148$ in an instrumented experiment with
17 million signatures.  The public response $z_1$ then contains an empirical
first-moment signal correlated with the two ternary polynomials $(s_0,s_1)$.

We give a streaming attack that uses fields decoded from ordinary valid
signatures.  Negacyclic normal equations estimate all 1,024 coefficients of
$(s_0,s_1)$, and a three-class mixture model gives an initial ternary guess.
The remaining ambiguity is completed with the public-key equation using a
binary-choice Kannan embedding.  Its parity structure makes $\{-1,1\}$ the
shortest possible magnitude for the centered correction coordinates; the
decoder explicitly rejects vectors without that shape.  Flatter followed by
BKZ-20 and BKZ-30 recovers the target vector.  We separate the exact score
identity from its heuristic scalar-leakage approximation and specify a finite
classifier and support search that can be fixed in advance for future tests.

On one deterministic DARTS-128 target, the experiment recovered all 1,536
generated ternary coefficients from 17,000,000 signatures.  Collection took
13.14 hours at 359.45 accepted signatures per second on a four-core Intel
Core i5-6500 with 16 GiB of RAM.  The original completion used retrospectively
tuned classifier and support parameters and took 4 minutes 55 seconds with
less than 87 MiB of resident memory.  A target-secret-independent execution
of the finite search on the same transcript recovered the same secret in
about 5 minutes.  Because that search was itself designed during work on this
target, the result does not estimate success over random key generation.  The
original secret seed $K$ was not recovered; packing the exact ternary secret
with a fresh $K$ constructed an equivalent signing key.  It produced a
1,449-byte signature on a message absent from the query set, and the
unmodified submitted verifier accepted it.  This is a successful fresh-message
forgery for one fixed DARTS-128 key.  No additional keys or parameter sets
have undergone end-to-end recovery tests.
\end{abstract}

\section{Introduction}

Fiat--Shamir-with-aborts signatures publish a response obtained by adding a
secret-dependent challenge product to a random mask.  Rejection sampling must
remove the dependence on the signing secret.  If one bimodal branch is
accepted more often than the other, repeated signatures can retain a small
multiple of the long-term key.

DARTS is an NGCC lattice-signature submission based on module-LWE and
module-SIS assumptions \cite{DARTS}.  It combines bimodal rejection sampling
\cite{DucasEtAl} with uniform hyperball sampling and a double-sphere output
distribution.  The response is compressed to obtain signatures of 1,449
bytes at the 128-bit level.  DARTS adds a compression-stability filter before
its sphere rejection.  That filter is branch dependent: it tests a shifted
commitment for one value of the bimodal bit and performs a tautological
comparison for the other.

The asymmetry is visible before any ternary-secret-recovery step is applied.
Under a uniform-commitment and independent-coordinate model, the gate alone
predicts that 44.133 percent of gate-passing samples use branch zero.  We
observe 44.148 percent among fully accepted signatures.  The numerical
agreement does not prove that later rejection is branch neutral.  In the
measured transcript, multiplying the public response by the adjoint of the
public challenge convolution and solving normal equations produces scores
correlated with both secret polynomials.

The estimator is not accurate enough to round every coefficient directly.
At 17 million signatures, a public mixture classifier leaves 47 wrong values
among 1,024 coefficients.  DARTS's public key supplies a much stronger final
test.  For DARTS-128 it satisfies
\[
             A_0s_0+A_1s_1+e=h\pmod q,
             \qquad s_0,s_1,e\in\{-1,0,1\}^{512},
\]
where $h$ has constant coefficient $(q+1)/2$ and all other coefficients zero.
Thus a candidate for $(s_0,s_1)$ determines $e$, and all 512 coefficients of
$e$ must be ternary.

Our contributions are:

\begin{itemize}[leftmargin=1.5em]
  \item We identify the branch-dependent compression gate in the specified
        DARTS signing algorithm and measure the final branch imbalance on one
        instrumented target.
  \item We derive an exact decomposition of a public normal-equation score
        for the two transmitted secret components and separate that identity
        from the heuristic scalar-leakage model.  The estimator is streaming
        and stores three ring accumulators.
  \item We disclose the retrospective choices used in the original
        130-variable completion and specify deterministic classifier-width
        and support ladders with a public validation rule.
  \item We state the precise decoding condition for the binary-choice Kannan
        embedding; shortness favors binary coordinates, while an explicit
        decoder and the complete public relation validate candidates.
  \item The experiment demonstrates exact ternary-secret recovery,
        equivalent signing-key construction, and an accepted fresh-message
        forgery from a 17-million-signature transcript on an ordinary
        four-core machine.
\end{itemize}

The attack uses no timing measurement, fault, malformed signature, weak hash
output, or exposed signing randomness.  The branch-dependent gate appears in
the specified procedure and is faithfully implemented by the reference code.

\section{Relevant parts of DARTS-128}
\label{sec:darts}

Let
\[
       \R_q=\Z_q[X]/(X^{512}+1), \qquad q=130817.
\]
DARTS-128 uses $(k,\ell)=(1,2)$.  Key generation samples three ternary
polynomials
\[
                    s=(s_0,s_1,e)\in\R_q^3
\]
with independent coefficient probabilities
\[
  \Prb[s_i=-1]=0.15,\quad
  \Prb[s_i=0]=0.70,\quad
  \Prb[s_i=1]=0.15.
\]
The public matrix is $A=(A_0,A_1,1)$ and the key relation is
\begin{equation}
                  A_0s_0+A_1s_1+e=h\pmod q,
\label{eq:public-relation}
\end{equation}
where $h=((q+1)/2,0,\ldots,0)$.  The public key determines $A_0$ and $A_1$.

For each signing attempt, the signer samples a hyperball mask $y$, a bimodal
bit $b$, and an auxiliary rejection bit.  It computes a commitment $w$, then
a sparse challenge $c$ of Hamming weight $\tau=30$.  The secret-dependent
response is
\begin{equation}
                     z=y+(-1)^bcs.
\label{eq:response}
\end{equation}
The signature directly encodes the two-polynomial block $z_1$ corresponding
to $(s_0,s_1)$, split into compressed high and low parts.  Anyone can decode
those fields and recover the integer coefficient vector used by verification.
The third response component is handled by a hint and is unnecessary for our
estimator.

Before computing Equation~\eqref{eq:response}, DARTS Algorithm~7 applies
\begin{equation}
           \Comp_d(w)=\Comp_d\bigl(w+(1-b)c\alpha\bigr),
\label{eq:gate}
\end{equation}
with $d=10$.  When $b=1$, Equation~\eqref{eq:gate} is
$\Comp_d(w)=\Comp_d(w)$ and always passes.  When $b=0$, adding the sparse
challenge can cross a compression boundary and force a restart.  The later
double-sphere rejection is intended to shape $z$.  The final accepted
transcript remained imbalanced in the instrumented experiment reported below.

\section{From the compression gate to a secret estimator}
\label{sec:leakage}

\subsection{Gate-only prediction and fully accepted measurement}

For a uniformly distributed coefficient $a\in\Z_q$, the gate model gives
\[
  \Prb[\Comp_d(a)=\Comp_d(a+1)]\simeq 1-\frac{2^d}{q}.
\]
Treating the $\tau$ challenged coefficients as independent gives the
heuristic branch-zero gate-pass probability
\begin{equation}
       \rho=\left(1-\frac{2^{10}}{130817}\right)^{30}
            \simeq 0.78997.
\label{eq:rho}
\end{equation}
The two branches begin with equal probability, so conditioning on this gate
model predicts
\begin{equation}
  \Prb[b=0\mid\text{gate passes}]
       =\frac{\rho}{1+\rho}\simeq0.44133,
  \qquad
  \E[(-1)^b]\simeq-0.11734.
\label{eq:branch-bias}
\end{equation}

This calculation covers the compression gate.  It does not model the later
sphere rejection, encoding checks, or dependencies among accepted variables.
In the 17-million-signature experiment, the hidden branch counter was read
only after data collection as a diagnostic.  It contained 7,505,081
branch-zero and 9,494,919 branch-one fully accepted signatures.  Thus the
observed values were
\[
       \widehat{\Prb}[b=0]=0.44147535,
       \qquad \bar\varepsilon=-0.11704929,
\]
numerically close to Equation~\eqref{eq:branch-bias}.  This agreement is an
empirical observation, not a proof that the later filters preserve the
gate-only distribution.  The attack itself does not use the hidden branch bit
or these counters.

\subsection{Public negacyclic normal equations}

Let $C_t$ be the $512\times512$ signed negacyclic convolution matrix of the
public challenge $c_t$ in signature $t$.  For response polynomial $p\in
\{0,1\}$, Equation~\eqref{eq:response} becomes
\begin{equation}
              z_{p,t}=y_{p,t}+\varepsilon_t C_t s_p,
              \qquad \varepsilon_t=(-1)^{b_t}.
\label{eq:linear-response}
\end{equation}
Define the streaming accumulators
\begin{equation}
       G=\sum_{t=1}^{N} C_t^TC_t,
       \qquad
       r_p=\sum_{t=1}^{N} C_t^Tz_{p,t}.
\label{eq:normal-acc}
\end{equation}
Negacyclic symmetry means $G$ is determined by one 512-coefficient
polynomial, so the complete attack state is $G,r_0,r_1$ plus counters and the
public key.  No raw signature storage is required.  After collection, solve
\begin{equation}
                         G\theta_p=r_p.
\label{eq:normal-solve}
\end{equation}

For a fixed accepted transcript, define
\[
  \eta_p=\sum_{t=1}^{N}C_t^Ty_{p,t},
  \qquad H=\sum_{t=1}^{N}\varepsilon_t C_t^TC_t.
\]
Then the following decomposition is exact:
\begin{equation}
             r_p=\eta_p+Hs_p,
       \qquad \theta_p=G^{-1}\eta_p+G^{-1}Hs_p,
\label{eq:exact-score}
\end{equation}
whenever $G$ is invertible.  This algebraic identity is the rigorous part of
the estimator derivation.

The often useful approximation
\begin{equation}
                       \theta_p\simeq\E[\varepsilon]s_p
\label{eq:scalar-heuristic}
\end{equation}
requires two additional concentration statements: $G^{-1}\eta_p$ must be
small, and $G^{-1}H$ must be close to the scalar matrix
$\E[\varepsilon]I$.  Those statements would follow from suitable
independence and zero-correlation hypotheses before conditioning.  They are
not proved for the distribution after the compression gate, sphere
rejection, and encoding checks.  Equation~\eqref{eq:scalar-heuristic} is
therefore a heuristic explanation, not a theorem about DARTS.

The target measurements also show that the scalar approximation is not
quantitatively exact.  Regressing the 17-million-signature scores on the
known secret only as a post hoc diagnostic gives slopes $-0.04160$ for
$s_0$ and $-0.04159$ for $s_1$, whereas the measured branch mean is
$\bar\varepsilon=-0.11705$.  The corresponding correlations are $-0.8563$
and $-0.9575$.  Thus later conditioning attenuates or otherwise modifies the
simple branch mean.  The first score also has a stable coordinate-dependent
component.  These effects prevent direct rounding but leave three empirical
class centers.

This estimator uses only $c_t$ and the decoded public $z_{1,t}$.  The chosen
messages may be arbitrary distinct messages.  No internal rejection count,
branch bit, mask, or secret coefficient appears in
Equations~\eqref{eq:normal-acc}--\eqref{eq:normal-solve}.

\subsection{Oracle branch-balancing diagnostic}

To test the role of branch proportions, we made one separate diagnostic run
of 100,000 fully accepted signatures on the same deterministic key.  The run
split the public normal equations by the instrumented hidden bit $b$.  An
oracle-balanced score then weighted each branch by the reciprocal of its
accepted count.  This diagnostic uses information unavailable to the attack;
it is a causality check rather than an attack step.

The run contained 44,072 branch-zero and 55,928 branch-one signatures.
Table~\ref{tab:balance} reports an ordinary least-squares slope in
$\theta=\alpha+\beta s$ and the coefficient correlation, pooling $s_0$ and
$s_1$.

\begin{table}[ht]
\centering
\caption{Same-key 100,000-signature oracle branch-balancing diagnostic.}
\label{tab:balance}
\begin{tabular}{lrr}
\toprule
score construction & slope $\beta$ & correlation\\
\midrule
branch zero only & $+0.36275$ & $+0.83206$\\
branch one only  & $-0.33965$ & $-0.83645$\\
observed mixture & $-0.03004$ & $-0.17900$\\
equal branch weight & $+0.01160$ & $+0.06989$\\
\bottomrule
\end{tabular}
\end{table}

Equal weighting removes most of the scalar signal and reverses its small
remainder in this sample; it does not produce an exact numerical zero.  A
100,000-signature diagnostic with dependent ring coefficients is insufficient
to distinguish finite-sample fluctuation from a smaller systematic residual.
The result supports branch imbalance as a major contributor, while stopping
short of claiming that balancing alone proves transcript independence.

\section{Public classification and support selection}
\label{sec:classification}

\subsection{Three-class mixture model and retrospective disclosure}
\label{sec:retrospective}

The original successful completion used classifier widths and support
parameters selected during exploratory work that included comparisons with
the target secret.  Once those parameters were fixed, the reported completion
used only the transcript and public key and applied a public validation rule.
We report it for reproducibility, but it is not an out-of-sample evaluation.
Sections~\ref{sec:public-classifier} and~\ref{sec:public-support} give a fully
specified target-secret-independent execution.

The coefficients of each $\theta_p$ are classified with the known ternary
prior $\pi=(0.15,0.70,0.15)$.  For class $a\in\{-1,0,1\}$, define
\begin{equation}
  L_{p,i}(a)=\log\pi_a-\log\sigma_p
      -\frac{(\theta_{p,i}-\mu_{p,a})^2}{2\sigma_p^2}.
\label{eq:likelihood}
\end{equation}
The original completion uses $\sigma_0=0.0140$ and $\sigma_1=0.0083$, with
initial ordered means $(0.0376,-0.0041,-0.0442)$.  On the target transcript,
100 expectation-maximization iterations update the three means using only
$\theta_p$ while retaining the public priors and fixed widths.  The final
guess and confidence are
\begin{align}
 \widehat s_{p,i}&=\arg\max_a L_{p,i}(a),\\
 \kappa_{p,i}&=\max_a L_{p,i}(a)
                 -\mathop{\rm secondmax}_a L_{p,i}(a).
\end{align}

The fitted target means were
\begin{center}
\begin{tabular}{c@{\qquad}rrr}
\toprule
polynomial & $a=-1$ & $a=0$ & $a=1$\\
\midrule
$s_0$ & $0.03589272$ & $-0.00592390$ & $-0.04756463$\\
$s_1$ & $0.04202248$ & $-0.00018468$ & $-0.04113097$\\
\bottomrule
\end{tabular}
\end{center}
Post hoc comparison with the target key found that the public guess was
correct in 977 of 1,024 positions: 45 errors were in $s_0$ and two were in
$s_1$.

\subsection{Deterministic public classifier search}
\label{sec:public-classifier}

A target-secret-independent execution must also fix how the classifier
parameters are obtained.  The following deterministic procedure uses only
the score and the public ternary prior.  For each polynomial, sort its 512
scores.  Initialize the mean for class $-1$ from the largest 77 values, the
mean for class 0 from the middle 358, and the mean for class $+1$ from the
smallest 77.  This ordering follows the negative leakage slope.  Initialize a
common component width with the public sample standard deviation, then run EM
with fixed priors while updating both the three means and the shared width:
\begin{align}
 \gamma_{i,a}&=\frac{\pi_a
  \exp\!\left(-({\theta_i-\mu_a})^2/(2\sigma^2)\right)}
  {\sum_{a'}\pi_{a'}
  \exp\!\left(-({\theta_i-\mu_{a'}})^2/(2\sigma^2)\right)},\\
 \mu_a&=\frac{\sum_i\gamma_{i,a}\theta_i}{\sum_i\gamma_{i,a}},
 \qquad
 \sigma^2=\frac1{512}\sum_{i,a}\gamma_{i,a}(\theta_i-\mu_a)^2.
\label{eq:public-em}
\end{align}
Stop when the largest mean or width change is below $10^{-13}$, or after 500
iterations.  Denote the resulting widths by
$\widehat\sigma_0,\widehat\sigma_1$.

To avoid choosing a width after seeing the target secret, enumerate the fixed
multiplier set
\[
             \Lambda=\{1,9/8,5/4,11/8,3/2\}.
\]
For every pair $(\lambda_0,\lambda_1)\in\Lambda^2$, run 100 mean-only EM
updates with widths $\lambda_p\widehat\sigma_p$, then run the public support
ladder below.  Order the 25 pairs first by
$(\max(\lambda_0,\lambda_1),\lambda_0+\lambda_1,\lambda_0,\lambda_1)$ and
break all remaining ties by index.  Traverse support rungs 0 through 5 within
each width pair.  At each trial, apply the fixed reduction schedule in
Section~\ref{sec:lattice}.  Accept only a candidate for which
$(s_0,s_1)$ is ternary, the publicly derived $e$ is ternary, all 512 public
equations hold, and a fresh signature verifies.  Otherwise advance to the
next support rung or width pair; if the finite grid is exhausted, report
failure at the fixed signature budget.

This defines a proper public attack experiment: fix the signature
budget, classifier search, support ladder, reduction schedule, and public stop
condition before drawing a target key.  The signing-oracle transcript and
public key are its only per-target inputs.  The procedure can be fixed in
advance for future tests.  One such evaluation could use a budget of 17
million signatures and record success or failure without reading the
generated secret.  We do not perform that additional-key experiment here.

\subsection{Exploratory support and a deterministic public ladder}
\label{sec:public-support}

Confidence alone misses some high-confidence $s_0$ errors caused by the
fixed component.  Their location is also visible publicly as blockwise drift.
Partition $\theta_0$ into sixteen consecutive blocks of 32 coefficients.  Let
$m_j$ be the mean of block $j$ and let $m_*$ be the median of the sixteen
block means.  Select the three blocks maximizing $|m_j-m_*|$.

At 17 million signatures, this rule ranks blocks 15, 0, and 1 first.  Their
means were respectively
\[
       -0.03754359,\quad -0.02454275,\quad -0.01962178,
       \qquad m_*=-0.00582224.
\]
The original successful completion took the union of these 96 block positions
and the 64 least-confident $s_0$ positions, leaving 128 positions after
overlap.  It added the two least-confident $s_1$ positions and used 48 public
equations, for 130 binary choices in a 179-dimensional lattice.

At each selected position, retain the most likely class and its runner-up.
Their difference is $a_j\in\{-1,1\}$.  Introduce one bit $x_j$ that chooses
between them.  Post hoc diagnostics show that all 47 wrong guesses lie in the
support and that the runner-up direction is correct for every one.  The final
reduction used only public data, and success was detected through
Equation~\eqref{eq:public-relation}.  These counts are the retrospectively
selected parameters disclosed in Section~\ref{sec:retrospective}.

For target-secret-independent execution, use the public ladder in
Table~\ref{tab:ladder}.  At each rung, select the stated number of
least-confident positions in each polynomial; unite the $s_0$ set with every
position in the stated number of blocks having largest $|m_j-m_*|$; and keep
the public runner-up class.  Ties are broken by increasing coefficient or
block index; likelihood ties are sorted by increasing class label and the
last class is preferred.  Construct the embedding with the stated prefix
length $m$ and stop only when all 1,536 derived coefficients are ternary and
the complete public relation holds.

\begin{table}[ht]
\centering
\caption{Deterministic public support ladder.}
\label{tab:ladder}
\begin{tabular}{crrrr}
\toprule
rung & low-confidence $s_0$ & low-confidence $s_1$
     & anomalous $s_0$ blocks & equations $m$\\
\midrule
0 & 32  & 2  & 1 & 32\\
1 & 64  & 2  & 2 & 40\\
2 & 64  & 2  & 3 & 48\\
3 & 96  & 8  & 4 & 56\\
4 & 128 & 16 & 6 & 64\\
5 & 160 & 32 & 8 & 72\\
\bottomrule
\end{tabular}
\end{table}

On the saved transcript, rung 2 publicly selects blocks 15, 0, and 1.  Its
unions contain 128 positions in $s_0$ and two in $s_1$, recreating the
130-choice, 179-dimensional instance.  The reduction schedule in
Section~\ref{sec:lattice} returns the same publicly validated ternary secret.

We also executed the deterministic public classifier procedure on the saved
transcript.  Its fitted base widths were
$(0.01037953,0.00715656)$.  Width-pair rung 2,
$(\lambda_0,\lambda_1)=(9/8,1)$, and support rung 2 produced a public guess
with 971 of 1,024 coefficients correct in a post hoc comparison.  The public
unions contained 127 positions in $s_0$ and two in $s_1$, giving a
178-dimensional lattice with 129 choices and 48 equations.  Flatter,
BKZ-20, and BKZ-30 returned a publicly valid vector selecting all 53 needed
corrections in 302.91 seconds of measured reduction time and 89,292 KiB peak
resident memory.  It again achieved exact ternary-secret recovery.  A basis
construction without access to the target secret produced a byte-identical
raw basis to the diagnostic build.

The ladder and width grid were designed after exploratory work on this target;
the implication for the strength of the claim is stated in
Section~\ref{sec:scope}.

\section{Binary lattice completion}
\label{sec:lattice}

\subsection{The modular correction problem}

Index the 130 selected corrections by $j$.  Let $D_j\in\Z_q^{m}$ be the first
$m=48$ coefficients of the public product obtained by placing $a_j$ at its
selected secret coordinate and zeros elsewhere.  Set
\begin{equation}
  D=(D_1|\cdots|D_u)\in\Z_q^{m\times u},\qquad u=130,
\end{equation}
and let
\begin{equation}
  r=\bigl(h-A_0\widehat s_0-A_1\widehat s_1\bigr)_{0:m-1}
       \pmod q.
\end{equation}
The desired correction bits satisfy
\begin{equation}
                         Dx+e'=r\pmod q,
       \qquad x\in\{0,1\}^{u},\quad e'\in\{-1,0,1\}^{m}.
\label{eq:binary-system}
\end{equation}

Ordinary integer embeddings admit many short vectors whose correction
coordinates are small integers but not valid class choices.  We instead
center each bit as $1-2x_j\in\{-1,1\}$.

\subsection{The embedding}

Use the following row basis of dimension $d=u+m+1$:
\begin{equation}
B=\begin{pmatrix}
 0_{m\times u} & 2qI_m & 0\\
 2I_u          & 2D^T   & 0\\
 \mathbf{1}_u^T& 2r^T   & 1
\end{pmatrix}.
\label{eq:basis}
\end{equation}

\begin{lemma}
If $(x,e')$ satisfies Equation~\eqref{eq:binary-system}, then the lattice
generated by Equation~\eqref{eq:basis} contains
\[
                    v=(\mathbf{1}-2x,\;2e',\;1).
\]
Conversely, any lattice vector with last coordinate $\pm1$ and first $u$
coordinates all in $\{-1,1\}$ can be oriented to last coordinate $1$ and
written $(\mathbf{1}-2x,2e',1)$ for some
$x\in\{0,1\}^{u}$ and $e'\in\Z^m$ satisfying $Dx+e'=r\pmod q$.
It solves Equation~\eqref{eq:binary-system} exactly when
$e'\in\{-1,0,1\}^m$.
\end{lemma}

\begin{proof}
Choose $k\in\Z^m$ such that $r-Dx+qk=e'$.  Add the final row of $B$, subtract
$x_j$ times binary row $j$, and add the corresponding modulus rows.  The
three blocks are $\mathbf{1}-2x$, $2e'$, and 1.  For the converse, orient the
vector so its last entry is 1.  Its integer combination contains the final
row once.  The first block is therefore $\mathbf{1}+2t$ for an integer vector
$t$; the assumed $\{-1,1\}$ shape gives
$x=(\mathbf{1}-v_{0:u})/2=-t\in\{0,1\}^u$.  Dividing the middle block by two
defines an integer $e'$ and gives $Dx+e'=r\pmod q$.  The ternary bound on
$e'$, the remaining 464 public equations, and the bounds on the complete
derived ternary vector
reject false candidates.
\end{proof}

The embedding does not force its first coordinates to be binary.  With last
coordinate oriented to 1, it forces only odd parity, so coordinates such as
$\pm3$ are also possible.  Euclidean shortness favors the minimum absolute
value 1, and the decoder explicitly accepts only the $\{-1,1\}$ shape.  The
lemma is an algebraic encoding and decoding statement; it is not a uniqueness
or shortest-vector theorem.  Recovery by the stated reduction schedule is an
experimental fact for the demonstrated instance.

On the demonstrated target, 12 of the first 48 coefficients of $e$ are
nonzero.  Hence
\begin{equation}
                    \norm{v}^2=130+4\cdot12+1=179.
\end{equation}
The lattice also has dimension 179 and determinant
\[
                         \det(B)=2^{130}(2q)^{48}.
\]
Its Gaussian-heuristic length is approximately 151.92, whereas the target
length is $\sqrt{179}\simeq13.38$, a ratio of 11.35.  This large separation
makes modest lattice reduction sufficient despite the structured basis.

\subsection{Reduction, validation, and signing}

The basis was first reduced with the iterated-compression method of Ryan and
Heninger \cite{RyanHeninger}, targeting root-Hermite factor 1.01.  We then
applied two loops of BKZ-20 followed by two loops of BKZ-30 with 256-bit MPFR
orthogonalization using fpylll 0.6.4 \cite{SchnorrEuchner,fpylll}.  The first
row after BKZ-30 had the required binary shape.  It selected 47 corrections.

The corrected $(s_0,s_1)$ determines
\begin{equation}
            e=h-A_0s_0-A_1s_1\pmod q,
\label{eq:derive-e}
\end{equation}
with centered representatives.  Every coefficient of $s_0,s_1,e$ was
ternary, so all 512 public equations validated.  Post hoc comparison showed
that the recovered 1,536 ternary coefficients equal those generated in the
original target key exactly.

DARTS stores the public key, $(s_0,s_1,e)$, and a secret seed $K$ in its
signing key.  The experiment does not recover the original $K$.  Once the
ternary relation is recovered, the attacker can choose a fresh $K$ and pack
an equivalent signing key.  The new seed changes only the deterministic
masks; it does not change Equation~\eqref{eq:public-relation}.  We therefore
describe the result precisely as \emph{exact ternary-secret recovery followed
by equivalent signing-key construction}.  The ordinary signing algorithm can
then sign fresh messages under the victim public key.

\section{End-to-end experiment}
\label{sec:experiment}

\subsection{Setup and controls}

The experiment used the submitted DARTS archive with SHA-256 digest
\begin{center}
\small\texttt{1846cfe63f0cef83e2e0ca21f5dcadce3c4b16da33713957be6d156e2a9e6e95}.
\end{center}
The machine had a four-core Intel Core i5-6500 at 3.20 GHz and 16 GiB of RAM,
running 64-bit Linux.  Data collection used eight OpenMP workers.  The
collector cached only key expansion and NTT values that are invariant across
messages; all sampling, hashing, rejection, rounding, compression, and
encoding steps were unchanged.

Three controls checked that this optimization did not create the leakage:
\begin{itemize}[leftmargin=1.5em]
  \item Across four keys and 128 messages per key, cached and submitted
        signers produced byte-identical signatures, identical accepted branch
        bits, and identical rejection counts for all 512 cases.  Every
        signature verified.
  \item Serial and four-thread collection produced byte-identical results and
        rejection counts for 1,000 messages, with zero mismatches.
  \item A separate four-thread run generated and verified 10,000 signatures,
        with zero failures.
\end{itemize}
The target key was generated deterministically for reproducibility.  Query
messages were distinct 16-byte counter encodings.  The final forgery message
was a different ASCII string and did not occur in the query set.

\subsection{Collection and statistical recovery}

Table~\ref{tab:collection} reports the collection result.  The wall time sums
three resumable segments.  Hidden branch counts and target-key comparisons
are diagnostic values read only after the public attack state was complete.

\begin{table}[ht]
\centering
\caption{DARTS-128 transcript collection on one target key.}
\label{tab:collection}
\begin{tabular}{lr}
\toprule
quantity & result\\
\midrule
accepted signatures & 17,000,000\\
signing attempts & 82,136,603\\
attempts per accepted signature & 4.83156\\
branch-zero / branch-one (diagnostic) & 7,505,081 / 9,494,919\\
total collection wall time & 47,294.659 s (13.137 h)\\
accepted-signature rate & 359.45 s$^{-1}$\\
collector peak resident memory & 3.3 MiB\\
public initial guess & 977 / 1,024 correct\\
errors in $s_0$ / $s_1$ (diagnostic) & 45 / 2\\
\bottomrule
\end{tabular}
\end{table}

Storing all fixed-size signatures would require about 24.6 GB.  The streaming
normal equations instead require only three 512-entry signed accumulators and
the public key.  The resulting public score state, excluding the integer
accumulators used for resumption, is 9,336 bytes.

\subsection{Completion and forgery}

Table~\ref{tab:completion} gives the original fixed-parameter completion
timings.
The signature collector was still occupying the same four-core machine, so
wall time includes substantial CPU contention.  Once its parameters were
fixed, the lattice run received no target secret, branch bits, or hidden
rejection data.  Section~\ref{sec:retrospective} discloses how those parameters
were selected.

\begin{table}[ht]
\centering
\caption{Original fixed-parameter completion timings.}
\label{tab:completion}
\begin{tabular}{lrr}
\toprule
stage & wall time & peak resident memory\\
\midrule
basis construction and public support & $<1$ s & 80 MiB\\
iterated-compression reduction & 41.80 s & 80.9 MiB\\
BKZ-20, two loops & 82.46 s & 81.7 MiB\\
BKZ-30, two loops & 169.71 s & 86.1 MiB\\
\midrule
streamlined completion total & approximately 295 s & 86.1 MiB\\
fresh signing and verification & 0.01 s & 1.8 MiB\\
\bottomrule
\end{tabular}
\end{table}

BKZ-30 returned a valid binary vector as its first basis row.  It changed
exactly 47 classified coefficients.  Equation~\eqref{eq:derive-e} produced a
fully ternary error polynomial and the complete public relation held.  After
the run, comparison with the known test key gave 1,024 of 1,024 correct
coefficients in $(s_0,s_1)$ and 512 of 512 in $e$.

We packed the recovered ternary secret with an independently chosen 32-byte
signing seed
and signed the fresh message
\begin{center}
\small\texttt{fresh DARTS message signed with the lattice-recovered equivalent key}.
\end{center}
The resulting signature had the specified length of 1,449 bytes.  The
unmodified submitted DARTS-128 verifier returned acceptance.  Thus the output
is a fresh-message forgery, rather than only a distinguisher or partial-key
leak.

\section{Scope, limitations, and mitigation}
\label{sec:scope}

The experiment establishes exact ternary-secret recovery, equivalent signing-
key construction, and a fresh-message forgery for one fixed DARTS-128 key.
In the formal EUF-CMA experiment, the adversary's success probability includes
the randomness of $\KeyGen$.  One successful key, particularly one used during
parameter exploration, does not estimate that probability.  The justified
claim is therefore \emph{a successful fresh-message forgery for one fixed
DARTS-128 key}.  Calling DARTS-128 practically broken in general would require
evidence that the method succeeds with non-negligible probability over fresh
keys.  Because the original seed $K$ was not recovered, we use the terms
\emph{exact ternary-secret recovery} and \emph{equivalent signing-key
construction} throughout.

Shorter profiling experiments on independently generated keys showed similar
three-class first-moment structure, but no additional end-to-end recovery was
attempted.  The mixture widths, original support parameters, public ladder,
and 17-million query count are empirical and were not fixed before the
reported experiment.  Their portability to fresh random keys is unmeasured.

We make no experimental claim for DARTS-256 or DARTS-512.  Their dimensions,
challenge weights, and compression parameters differ.  The same structural
question applies because Algorithm~7 uses the branch-dependent factor
$(1-b)$ at every level, but practical extrapolation requires separate data.

The attack is not an implementation-only side channel.  Equation~\eqref{eq:gate}
appears in the normative signing algorithm, and the reference implementation
faithfully performs it.  Restricting timing information or rewriting the code
in constant time does not remove the public transcript bias.

A candidate repair should make the commitment filter independent of the
bimodal branch, or use a symmetric condition whose accepted response
distribution has a proof covering the complete compressed transcript.  The
oracle diagnostic shows that equal branch weighting removes most of the
measured scalar signal, but it does not prove that balancing alone eliminates
all public dependence.  Any revision should be tested for public correlations
$\sum C_t^Tz_t$ at large query counts and supported by a new distribution
argument.

Seventeen million chosen-message queries are more than typical protocol use,
but they were practical for this instance: the complete transcript was
generated in about 13 hours on a decade-old four-core processor.  Query limits
can reduce exposure, while a cryptographic repair requires addressing the
accepted-transcript dependence itself.

\section{Public record and novelty check}

As of 26 September 2026, the public DARTS page at ngcc.dev listed one finding:
a 256-bit collision ceiling for the DARTS-512 message representative
\cite{NGCCDARTS}.  It did not describe branch-dependent leakage, exact
ternary-secret recovery, or a DARTS-128 forgery.  We also searched the public
NGCC coverage, the NIST pqc-forum index, and the Cryptology ePrint index.  The
contemporary overview of early NGCC findings likewise reports no DARTS
ternary-secret-recovery attack \cite{Ivezic}.  We found no prior public account
of this attack.  This novelty check is necessarily preliminary and awaits
human review.

\section{Conclusion}

DARTS's compression-stability gate treats its two bimodal branches
differently, and the fully accepted transcript was branch imbalanced in the
reported experiment.  Public responses from 17 million signatures yielded
scores correlated with the two ternary secret polynomials; a binary-choice
lattice embedding then used the public relation to repair the remaining
errors.  On the demonstrated key, collection took about 13.14 hours and the
original completion took about five minutes on an ordinary four-core machine.
The experiment recovered the exact ternary secret, constructed an equivalent
signing key with a fresh seed, and produced an accepted signature on a fresh
message.

\begin{thebibliography}{10}

\bibitem{DARTS}
Y.~Pan, Q.~Lu, Y.~Wu, J.~Li, C.~Cheng, R.~Liu, and Q.~Du.
\newblock \emph{DARTS: A Lattice-Based Digital Signature Scheme -- Algorithm
  Specification and Supporting Documentation}.
\newblock NGCC submission, 29 June 2026.

\bibitem{Lyubashevsky}
V.~Lyubashevsky.
\newblock Fiat--Shamir with aborts: Applications to lattice and
  factoring-based signatures.
\newblock In \emph{ASIACRYPT 2009}, pages 598--616, 2009.

\bibitem{DucasEtAl}
L.~Ducas, A.~Durmus, T.~Lepoint, and V.~Lyubashevsky.
\newblock Lattice signatures and bimodal Gaussians.
\newblock In \emph{CRYPTO 2013}, pages 40--56, 2013.

\bibitem{DeveveyEtAl}
J.~Devevey, O.~Fawzi, A.~Passelegue, and D.~Stehle.
\newblock On rejection sampling in Lyubashevsky's signature scheme.
\newblock In \emph{ASIACRYPT 2022}, pages 34--64, 2022.

\bibitem{RyanHeninger}
K.~Ryan and N.~Heninger.
\newblock Fast practical lattice reduction through iterated compression.
\newblock In \emph{CRYPTO 2023}, pages 3--36, 2023.

\bibitem{SchnorrEuchner}
C.-P. Schnorr and M.~Euchner.
\newblock Lattice basis reduction: Improved practical algorithms and solving
  subset sum problems.
\newblock \emph{Mathematical Programming}, 66(1):181--199, 1994.

\bibitem{fpylll}
The FPLLL development team.
\newblock \emph{fpylll, a Python wrapper for the fplll lattice reduction
  library}, version 0.6.4, 2025.
\newblock \url{https://github.com/fplll/fpylll}.

\bibitem{NGCCDARTS}
ngcc.dev.
\newblock DARTS (sign-08) public findings.
\newblock \url{https://ngcc.dev/reports/sign-08.html}, accessed 26 September
  2026.

\bibitem{Ivezic}
M.~Ivezic.
\newblock China's Next-Generation Crypto Candidates Drew 104 Public Findings
  in Three Days.
\newblock \url{https://postquantum.com/security-pqc/china-ngcc-candidate-flaws/},
  23 September 2026.

\end{thebibliography}

\end{document}
